% !TEX program = pdflatex
\documentclass[11pt,a4paper]{article}
\usepackage[margin=1.02in]{geometry}
\usepackage{amsmath,amssymb,amsthm,mathtools}
\usepackage{booktabs,microtype}
\usepackage{xcolor}
\usepackage[numbers,sort&compress]{natbib}
\usepackage[colorlinks=true,linkcolor=blue!55!black,citecolor=blue!55!black,urlcolor=blue!55!black]{hyperref}

\newtheorem{theorem}{Theorem}[section]
\newtheorem{proposition}[theorem]{Proposition}
\newtheorem{lemma}[theorem]{Lemma}
\newtheorem{corollary}[theorem]{Corollary}
\theoremstyle{definition}
\newtheorem{definition}[theorem]{Definition}
\theoremstyle{remark}
\newtheorem{remark}[theorem]{Remark}
\newcommand{\E}{\mathbb{E}}
\newcommand{\Pp}{\mathbb{P}}
\newcommand{\R}{\mathbb{R}}
\newcommand{\Ccal}{\mathcal{C}}
\newcommand{\norm}[1]{\lVert#1\rVert}

\title{\textbf{Half-Wormholes as Wiener Chaos:\ Correlated-Replica Spectroscopy, Exact Pfaffian Sectors, and Boundary Moment Indeterminacy}}
\author{}
\date{Working draft -- 28 October 2026}
\hypersetup{pdfauthor={}}

\begin{document}
\maketitle

\begin{abstract}
The fixed-coupling partition function of one-time-point SYK is a homogeneous multilinear
polynomial of degree $n=N/q$ in independent Gaussian couplings.  Its square therefore has
a canonical Wiener--It\^o chaos decomposition.  This is a probabilistic refinement of the
exact half-wormhole expansion already constructed by Mukhametzhanov, not a replacement of
that prior result.  We derive five consequences of the canonical grading.  First, two
replicas with coupling correlation $\rho$ obey
$\E[z(J)z(J')]=\rho^n\E[z^2]$; scanning $\rho$ diagonalises and measures the chaos-sector
norms of $z^2$.  Second, for $q=2$ we reduce the previously known general variance sum to the closed formula
$\E[W_{2k}^2]=(2n)!\binom{2(n-k)}{n-k}4^{-(n-k)}$.  It yields an explicit correlated
determinant formula for two Gaussian antisymmetric matrices.  Third, the same profile gives
$\E[z^4]/\E[z^2]^2=2n+1$ and quantifies the failure of a wormhole-plus-Wick-square
truncation at $q=2$.  Fourth, Gaussian hypercontractivity supplies non-asymptotic,
single-draw reconstruction bounds sector by sector.  Fifth, the boundary moment problem for
the law of $z$ is determinate for $n\le2$ and indeterminate for every $n\ge3$.  We prove the
latter by a conditional-Gaussian density lower bound and Krein's criterion.  The result is a
complete statement for the zero-dimensional model.  It does not establish a dynamical bulk
geometry or a complete theory of quantum gravity.
\end{abstract}

\section{Introduction}

In the one-time-point SYK model, fixed-coupling factorisation is represented by a finite
polynomial identity.  Saad, Shenker, Stanford and Yao identified the wormhole and
coupling-sensitive half-wormhole saddles \citep{SaadEtAl2024}.  Mukhametzhanov subsequently
constructed an exact finite expansion of $z^2$, including orthogonality, closed sector
variances and the normal-ordered-square interpretation of the leading half-wormhole
\citep{Mukhametzhanov2022}.  Related graded and Hermite structures appear in
\citet{TianYang2022}.  Those results are prior art and are used here as such.
After matching conventions, Mukhametzhanov's term $\Phi_{n-k}$ is precisely the chaos
projection $W_{2k}$ below.  Thus neither the finite sector decomposition nor its general
variance sum is claimed as new; the $q=2$ binomial form is a closed reduction supplied with
an independent generating-function proof.

Our starting point is narrower: over the Gaussian coupling ensemble, $z^2$ has a unique
decomposition into eigenspaces of the Ornstein--Uhlenbeck (OU) semigroup.  Identifying this
canonical grading makes the correlated-replica kernel, Gaussian hypercontractivity and the
Hamburger moment problem available without introducing new saddles. A targeted search
recorded in the July 2026 source project found no substantive duplicate of the combined
results in the half-wormhole, random-matrix and moment-problem literature. That historical
search is bounded and has not been replaced by a comprehensive current novelty audit for
this repository edition.

The exact contributions are:
\begin{enumerate}
\item the $\rho^n$ two-replica identity and OU sector spectroscopy;
\item a closed $q=2$ reduction of the known sector-variance sum, an independent enumeration,
and the correlated antisymmetric-determinant corollary;
\item an explicit fourth-moment and truncation diagnosis, together with single-draw bounds;
\item the sharp determinacy dichotomy at chaos degree $n=2/3$ and its spectral consequences.
\end{enumerate}
We separate exact finite-$N$ theorems from asymptotic interpretations.  In particular, the
standard Nualart--Peccati fourth-moment equivalence applies directly to sequences in a
\emph{fixed} Wiener chaos \citep{NualartPeccati2005,NourdinPeccati2012}.  In the conventional
SYK scaling with fixed $q$ and $N\to\infty$, the chaos order $n=N/q$ grows; fourth moments
alone are then not advertised here as a central-limit proof.  The $L^2$ truncation statements
we use instead follow from the exact sector variances.

\section{Model and canonical grading}

Let $q\ge2$ be even, $N=qn$, and let $A$ run over the $q$-subsets of $[N]$.  The couplings
$J_A$ are independent standard normal variables.  Berezin integration gives, up to an
irrelevant global phase,
\begin{equation}
 z(J)=\sum_{P\in\operatorname{Part}_q(N)}\epsilon(P)\prod_{A\in P}J_A,
 \label{eq:z}
\end{equation}
where $P$ partitions $[N]$ into $n$ blocks of size $q$ and $\epsilon(P)=\pm1$ is the
Grassmann sign.  For $q=2$, $z(J)=\operatorname{Pf}J$ for the antisymmetric matrix with
upper-triangular entries $J_{ij}$.

Let $L^2(\gamma)=\bigoplus_{d\ge0}\Ccal_d$ be the Wiener--It\^o decomposition of the
Gaussian coupling space \citep{Ito1951,Janson1997}.  The OU operator is
\begin{equation}
 (T_\rho F)(J)=\E_K F\!\left(\rho J+\sqrt{1-\rho^2}K\right),
 \qquad -1\le\rho\le1,
 \label{eq:OU}
\end{equation}
and acts on $\Ccal_d$ by multiplication by $\rho^d$.

\begin{proposition}[Purity and symmetry]\label{prop:purity}
$z\in\Ccal_n$, $\E[z^2]=|\operatorname{Part}_q(N)|$, and the law of $z$ is symmetric.
\end{proposition}

\begin{proof}
Every monomial in \eqref{eq:z} is a product of $n$ distinct first Hermite polynomials, hence
belongs to $\Ccal_n$.  Distinct partitions give distinct monomials, so orthogonality gives
the second statement.  Fix a vertex, say $1$, and change the sign of every coupling $J_A$
with $1\in A$.  Every partition contains exactly one such block, so $z\mapsto-z$, whereas
the joint Gaussian law is unchanged.  Thus all odd moments vanish and the law is symmetric.
\end{proof}

\begin{theorem}[Correlated-replica identity]\label{thm:replica}
Let $J'=\rho J+\sqrt{1-\rho^2}K$ with $K$ an independent coupling draw.  Then
\begin{equation}
 \E[z(J)z(J')]=\rho^n\E[z^2].
 \label{eq:replica}
\end{equation}
More generally, if $X=\sum_dX_d$ and $Y=\sum_dY_d$ are their chaos decompositions, then
\begin{equation}
 \E[X(J)Y(J')]=\sum_d\rho^d\E[X_dY_d].
 \label{eq:mehler}
\end{equation}
\end{theorem}

\begin{proof}
Equation \eqref{eq:mehler} is Mehler's identity
$\E[X(J)Y(J')]=\langle X,T_\rho Y\rangle$.  Proposition~\ref{prop:purity} then gives
\eqref{eq:replica}.
\end{proof}

Define the canonical half-wormhole sectors by
\begin{equation}
 z^2=\sum_{k=0}^{n}W_{2k},\qquad W_{2k}\in\Ccal_{2k},\qquad
 \E[W_{2k}W_{2\ell}]=0\quad(k\ne\ell).
 \label{eq:sectors}
\end{equation}
Thus $W_0=\E[z^2]$, while $W_{2n}=:\!z^2\!:$ is the Wick square.  This decomposition is
canonical because the $\Ccal_d$ are the distinct OU eigenspaces.  It need not coincide
term-by-term with an arbitrary orthogonal expansion once $n\ge3$.  For the normal-ordered
expansion of \citet{Mukhametzhanov2022}, however, homogeneity gives the exact identification
$\Phi_{n-k}=W_{2k}$ after matching the coupling normalisation.

Applying \eqref{eq:mehler} to $z^2$ gives
\begin{equation}
 \E[z^2(J)z^2(J')]=\sum_{k=0}^{n}\rho^{2k}\E[W_{2k}^2].
 \label{eq:spectroscopy}
\end{equation}
Hence correlation in coupling space is a sector spectrometer: the coefficients of the
polynomial in $\rho^2$ are precisely the sector norms.

\section{Exact Pfaffian sector spectrum}

For $q=2$, ordered pairs of perfect matchings provide a closed enumeration.

\begin{theorem}[$q=2$ sector profile]\label{thm:profile}
For $z=\operatorname{Pf}J$ with $J$ a $2n\times2n$ Gaussian antisymmetric matrix,
\begin{equation}
 \boxed{\quad\E[W_{2k}^2]=(2n)!\binom{2m}{m}4^{-m},\qquad m=n-k,\quad0\le k\le n.\quad}
 \label{eq:profile}
\end{equation}
Equivalently,
\begin{equation}
 \sum_{n\ge0}\sum_{k=0}^{n}\E[W_{2k}^2]x^{2k}\frac{y^{2n}}{(2n)!}
 =\frac{1}{(1-x^2y^2)\sqrt{1-y^2}}.
 \label{eq:GF}
\end{equation}
\end{theorem}

\begin{proof}
Expand $z^2$ over ordered perfect matchings $(P,Q)$.  Their symmetric difference is a
vertex-disjoint union $U$ of alternating even cycles of length at least four; shared edges
contribute $J_e^2=H_2(J_e)+1$.  A chaos monomial is therefore labelled by a set $T$ of
excited shared edges and by $U$, with order $2|T|+|E(U)|$.  Each alternating cycle has two
splittings into $(P,Q)$, and the unused vertices carry an arbitrary common matching.  Upon
squaring a coefficient, the two common matchings form another union of doubled edges and
alternating cycles.  Sign coherence holds because matching edges have even block size.

The squared norm is thus the exponential generating function of labelled components of four
types: an excited doubled edge of weight $2x^2$, a first-family cycle of length $2j\ge4$
of weight $4x^{2j}$, an unexcited doubled edge of weight $1$, and a second-family cycle of
length $2j\ge4$ of weight $2$.  Since the number of undirected alternating cycles on $2j$
labelled vertices supplies the standard $1/(2j)$ exponential weight, the exponential formula
gives
\[
 \exp\!\left[x^2y^2+\sum_{j\ge2}\frac{(xy)^{2j}}{j}
 +\frac{y^2}{2}+\sum_{j\ge2}\frac{y^{2j}}{2j}\right]
 =\frac{1}{(1-x^2y^2)\sqrt{1-y^2}}.
\]
Extracting $x^{2k}y^{2n}$ and using
$[u^m](1-u)^{-1/2}=\binom{2m}{m}4^{-m}$ proves \eqref{eq:profile}.
\end{proof}

Equation~\eqref{eq:profile} is the $q=2$ evaluation of the general finite variance sum in
\citet{Mukhametzhanov2022}, not a priority claim for the sector variances.  The proof above is
an independent labelled-component derivation that exposes the rational generating function
used in the following consequences.

Two checks are immediate.  At $k=0$, \eqref{eq:profile} equals
$\E[z^2]^2=((2n-1)!!)^2$.  Summing over $k$ and using
\begin{equation}
 \sum_{m=0}^{n}\binom{2m}{m}4^{-m}
 =(2n+1)\binom{2n}{n}4^{-n}
\end{equation}
gives
\begin{equation}
 \frac{\E[z^4]}{\E[z^2]^2}=2n+1.
 \label{eq:fourthq2}
\end{equation}

\begin{corollary}[Correlated antisymmetric determinants]\label{cor:det}
If $J'=\rho J+\sqrt{1-\rho^2}K$ with $J,K$ independent Gaussian antisymmetric
$2n\times2n$ matrices, then
\begin{equation}
 \E[\det J\,\det J']
 =(2n)!\sum_{m=0}^{n}\binom{2m}{m}4^{-m}\rho^{2(n-m)}.
 \label{eq:det}
\end{equation}
\end{corollary}

\begin{proof}
$\det J=\operatorname{Pf}(J)^2=z^2$.  Combine \eqref{eq:spectroscopy} with
Theorem~\ref{thm:profile}.
\end{proof}

Formula \eqref{eq:det} is distinct from the general characteristic-polynomial averages
treated in random-matrix theory \citep{deBruijn1955,MehtaNormand2001}.  We have not found
this correlated antisymmetric specialisation stated in the searched literature.  This
negative search is recorded as evidence, not as a mathematical proof of priority.

\section{Truncation and single-draw control}

At $q=2$, the fraction of the total $L^2$ norm in sector $k$ is proportional to
$\binom{2(n-k)}{n-k}4^{-(n-k)}$.  The Wick-square fraction is
\begin{equation}
 \frac{\E[W_{2n}^2]}{\E[z^4]}
 =\frac{1}{(2n+1)\binom{2n}{n}4^{-n}}
 \sim\frac{\sqrt{\pi}}{2\sqrt n}.
 \label{eq:topfrac}
\end{equation}
Thus no top sector dominates, and the wormhole-plus-Wick-square truncation fails in relative
$L^2$ as $n\to\infty$.  Equation \eqref{eq:fourthq2} also shows divergence of the
normalised fourth moment.  This rules out Gaussian convergence in fourth moment; we do not
infer a general weak-limit theorem from this fact alone.

For any sector, hypercontractivity converts its exact norm into a single-draw bound.

\begin{lemma}[Chaos tail envelope]\label{lem:tail}
If $X\in\Ccal_d$, then for $t>0$
\begin{equation}
 \Pp(|X|\ge t\norm{X}_2)
 \le B_d(t):=\inf_{p\ge2}(p-1)^{pd/2}t^{-p}.
 \label{eq:B}
\end{equation}
In particular $B_d(t)\le t^{-2}$, and for $t\ge(2e)^{d/2}$,
$B_d(t)\le\exp[-(d/2e)t^{2/d}]$.
\end{lemma}

\begin{proof}
Gaussian hypercontractivity gives
$\norm{X}_p\le(p-1)^{d/2}\norm{X}_2$ \citep{Nelson1973,Janson1997}.  Markov's inequality
gives \eqref{eq:B}.  The two displayed bounds follow from $p=2$ and from the admissible
choice $p=t^{2/d}/e$, respectively.
\end{proof}

Put $f_k=\E[W_{2k}^2]/\E[z^2]^2$ and
$E_{\rm tr}=z^2-\E[z^2]-W_{2n}=\sum_{k=1}^{n-1}W_{2k}$.

\begin{theorem}[Single-draw reconstruction]\label{thm:single}
For any $\varepsilon_k>0$ with $\sum_{k=1}^{n-1}\varepsilon_k=\varepsilon$,
\begin{equation}
 \Pp\!\left(|E_{\rm tr}|\ge\varepsilon\E[z^2]\right)
 \le\min\left\{\frac{\sum_{k=1}^{n-1}f_k}{\varepsilon^2},
 \sum_{k=1}^{n-1}B_{2k}\!\left(\frac{\varepsilon_k}{\sqrt{f_k}}\right)\right\}.
 \label{eq:single}
\end{equation}
\end{theorem}

\begin{proof}
Orthogonality gives $\E[E_{\rm tr}^2]/\E[z^2]^2=\sum_{k=1}^{n-1}f_k$, hence the first bound.
For the second, if $|\sum W_{2k}|\ge\sum\varepsilon_k\E[z^2]$, at least one sector exceeds
its threshold.  Apply Lemma~\ref{lem:tail} and the union bound.
\end{proof}

Theorem~\ref{thm:single} is non-asymptotic and valid for one fixed coupling draw.  It does
not assert that its dimension-free constants are optimal.  When prior exact sector estimates
show $\sum_{k=1}^{n-1}f_k\to0$ in a specified large-$N$ regime \citep{Mukhametzhanov2022}, the
first term already proves relative $L^2$ reconstruction; the tail term refines atypical
events sector by sector.

\section{Boundary moment determinacy}

Let $\mu$ be the law of $z$.  In the baby-universe/GNS dictionary, the moments
$\mu_j=\E[z^j]$ are multi-boundary data, and multiplication by $z$ acts on the closure of
polynomials in $L^2(\mu)$ \citep{MarolfMaxfield2020,JanssenEtAl2021}.  Whether all moments
determine $\mu$ is exactly the Hamburger moment problem.

\begin{theorem}[Determinacy dichotomy]\label{thm:moment}
For the one-time-point model \eqref{eq:z}:
\begin{enumerate}
\item if $n\le2$, the law $\mu$ is Hamburger determinate;
\item if $n\ge3$, the law $\mu$ is Hamburger indeterminate.
\end{enumerate}
\end{theorem}

\begin{proof}
For $n\le2$, hypercontractivity applied to $z\in\Ccal_n$ gives
\[
 \mu_{2m}^{1/(2m)}=\norm{z}_{2m}
 \le(2m-1)^{n/2}\norm{z}_2.
\]
Therefore $\sum_m\mu_{2m}^{-1/(2m)}$ dominates a constant multiple of
$\sum_m m^{-n/2}$ and diverges.  Carleman's criterion implies determinacy
\citep{Akhiezer1965,Simon1998}.

Now let $n\ge3$.  Expand along blocks containing vertex $1$:
\begin{equation}
 z=\sum_{A\ni1}\epsilon_AJ_A C_A,
 \label{eq:cofactor}
\end{equation}
where every cofactor $C_A$ uses only couplings on subsets not containing $1$ and is
independent of the family $\{J_A:A\ni1\}$.  Conditionally on the cofactors,
$z\sim\mathcal N(0,S^2)$ with $S^2=\sum_{A\ni1}C_A^2$.  Since the cofactor vector is not
identically zero and its common zero set is algebraic of Gaussian measure zero, $S>0$
almost surely.  Thus $\mu$ has a strictly positive Gaussian-mixture density
\begin{equation}
 f(t)=\E\left[(2\pi S^2)^{-1/2}\exp\!\left(-\frac{t^2}{2S^2}\right)\right].
 \label{eq:density}
\end{equation}

We need a lower tail bound on $f$.  Fix a partition
$\{A_0,B_2,\ldots,B_n\}$ with $1\in A_0$.  On the event that
$J_{B_j}\in[a,a+1]$ for $j=2,\ldots,n$ and all other couplings not containing $1$ have
absolute value at most one, the leading monomial of $C_{A_0}$ has degree $n-1$.  All its
other monomials, and every $C_A$ with $A\ne A_0$, contain at most $n-2$ of the selected
large variables.  Hence, for $a$ above a model-dependent constant,
\[
 c a^{n-1}\le S\le C a^{n-1},\qquad
 \Pp(E_a)\ge c_1e^{-c_2a^2}.
\]
Insert this event into \eqref{eq:density} and choose $a=t^{1/n}$.  For sufficiently large
$t$,
\begin{equation}
 f(t)\ge c_3t^{-(n-1)/n}\exp[-c_4t^{2/n}].
 \label{eq:lower}
\end{equation}
Because $2/n<1$, \eqref{eq:lower} implies
\[
 \int_{\R}\frac{-\log f(t)}{1+t^2}\,dt<\infty
\]
up to an irrelevant finite contribution on compact intervals.  Krein's criterion therefore
implies Hamburger indeterminacy \citep{Krein1945,Stoyanov2000}.  This proves the dichotomy.
\end{proof}

For $n\ge3$, the symmetric Jacobi operator associated with the boundary moments is not
essentially self-adjoint on polynomials; its canonical self-adjoint extensions form the
Nevanlinna family and have discrete extremal spectral measures \citep{Akhiezer1965,Simon1998}.
The physical coupling law \eqref{eq:density} is continuous and non-extremal, so polynomials
are not dense in its full $L^2$ space.  This does not make the GNS construction inconsistent.
It says only that one-time boundary moments do not select the microscopic coupling measure
uniquely once one boundary amplitude is a Gaussian polynomial of degree at least three.

\section{Scope and conclusion}

The exact half-wormhole expansion, its identification with the chaos projections, and its
general sector-variance sum are not new.  Our contribution is the correlated-replica use of
the OU grading and the additional consequences developed here.  Equations \eqref{eq:replica} and
\eqref{eq:spectroscopy} turn correlated replicas into a sector probe.  Theorem
\ref{thm:profile} gives a closed Pfaffian reduction and independent proof, while
Corollary~\ref{cor:det} gives the correlated-determinant consequence.  Theorem~\ref{thm:single} is a finite single-draw statement.
Theorem~\ref{thm:moment} sharply separates determinate and indeterminate boundary laws.

The limitations are equally sharp.  The model is zero-dimensional; ``wormhole'' labels the
ensemble-connected term, not a derived spacetime geometry.  Gaussianity is essential to the
OU and hypercontractive form used here.  A thermal partition function has an infinite chaos
expansion.  Finally, when the chaos order grows with $N$, the fixed-chaos fourth-moment
theorem cannot by itself prove a central limit.  These boundaries prevent the manuscript
from being read as a construction of complete quantum gravity; it is a rigorous result about
a solvable factorisation model.

\section*{AI-assisted workflow and unresolved attribution}
The following statement is supplied by the repository operator:
\begin{quote}
I genuinely do not know how authorship, intellectual credit, or ownership should be
attributed for these results. My workflow was to ask AI systems to conduct literature
research, autonomously select open problems, and attempt to solve them. The principal
research work was carried out mainly by Fable 5. DeepSeek and several GPT models
participated in cross-review, checking, and manuscript preparation and organization.
These manuscripts and their scientific claims have not yet undergone human expert review.
\end{quote}
This describes the overall workflow reported by the operator; a complete per-model,
per-manuscript allocation of contributions is unavailable. Author and affiliation fields
are intentionally blank. Automated checks and AI internal reviews are not human peer
review. This working draft makes no claim of journal acceptance or settled publication
priority. Historical source dates are recorded separately in the repository history.

\bibliographystyle{plainnat}
\bibliography{references}
\end{document}
